\subsection{Условная вероятность. Независимые события}
Давайте определим как лучше определить условную вероятность. Будем обозначать вероятность события $A$ при условии события $B$ как $\mathbb{P}(A \mid B)$. Понятно, что если мы предполагаем, что событие $B$ уже наступило, то событие $B$ уже становится пространством элементарных исходов, а все события из алгебры событий теперь ограничиваются элементарными исходами $A$, т.е. каждое событие $A$ заменяется на событие $AB$. Тем самым, довольно логичным определением выглядит следующее:

\begin{definition}[Условная вероятность]
Условной вероятностью события $A$ при $B$ называют вероятность наступления события $A$ при условии наступления события $B$ ($\mathbb{P}(B) \neq 0$), которая равна
\[
\mathbb{P}(A \mid B) = \frac{\mathbb{P}(AB)}{\mathbb{P}(A)}.
\]
\end{definition}

Обратим внимание на условие $\mathbb{P}(B) \neq 0$. Логично, что вероятность события $A$ при наступлении события $B$, если событие $B$ невозможно, бессмысленна. Но ведь нулевая вероятность не означает невозможность события.

Такую проблему мы решим в следующем разделе, когда изучим понятия интеграла Лебега и математического ожидания.

\begin{definition}[Независимые события]
События $A$ и $B$ называются \textbf{независимыми}, если вероятность наступления одного события не зависит от наступления другого:
\[
\mathbb{P}(A \mid B) = \mathbb{P}(A).
\]
\end{definition}

Заметим сразу, что если $\mathbb{P}(A \mid B) = \mathbb{P}(A)$, то выполняется формула умножения вероятностей.

\begin{theorem}[Формула умножения вероятностей]
Если $A$ и $B$ -- независимые события, то
\[
\mathbb{P}(AB) = \mathbb{P}(A) \cdot \mathbb{P}(B).
\]
\end{theorem}

\begin{definition}[Попарно независимые события]
События $\{A_i\}_{i \in I}$ называются попарно независимыми, если любая пара событий из этого набора независима.
\end{definition}

\begin{definition}[Независимые в совокупности события]
События $\{A_i\}_{i \in I}$ называются независимыми в совокупности, если для любого конечного подмножества $I' \subseteq I$ верно
\[
\mathbb{P}\!\left(\bigcap_{k=1}^{n} A_{i_k}\right) = \prod_{k=1}^{n} \mathbb{P}(A_{i_k}).
\]
\end{definition}

\begin{notice}
Очевидно, что из независимости в совокупности следует независимость попарная. Однако обратное неверно.
\end{notice}